\section{Cheminement méthodologique et résultats intermédiaires}
\label{app:path}
% =====================================================================

Cette annexe détaille dans l'ordre chronologique les étapes du Tableau~\ref{tab:evolution}.

\paragraph{1. Les fichiers de HOD ne sont pas des réalisations.} Une première version de la chaîne moyennait les statistiques de 64 fichiers de HOD différents par cosmologie pour construire des ensembles de « covariance » et de « dérivées », et les mesurait dans le cube central de $500\Mpch$ de chaque boîte. Les fichiers de HOD sont des modèles de HOD différents, pas des réalisations : leur dispersion ($\sim55\%$) dominait les différences cosmologiques ($\sim1\%$), et les différences finies étaient compatibles avec du bruit ($|S/N|<1$ médian par bin pour $\omega_b$, $n_s$, $\sigma_8$). Cette version présentait aussi cinq problèmes d'implémentation, tous corrigés : les comptages de paires supposaient des conditions aux limites périodiques dans une boîte découpée non périodique ; la 2PCF utilisait la ligne de visée au point médian depuis le centre de la boîte au lieu de l'axe selon lequel les RSD avaient été appliquées, ce qui effaçait le quadrupôle ; les quartiles étaient calculés conjointement sur les galaxies et les points aléatoires, donnant des effectifs inégaux et dépendant de la réalisation ; la triangulation de Delaunay était calculée sur la boîte découpée sans extension périodique ; et ASTRA était d'abord appliqué au seul catalogue aléatoire, de sorte que les quartiles aléatoires ne portaient aucune information sur le champ de galaxies et que leurs corrélations étaient compatibles avec zéro.

\paragraph{2. Paires appariées en HOD.} Nous avons ensuite sélectionné, pour chaque paire $\pm$, les deux catalogues aux HOD les plus proches (deux fois plus proches que des paires aléatoires, Tableau~\ref{tab:hods}) et les avons différenciés sous-volume par sous-volume (Section~\ref{sec:finitediff}). Les différences sont précises mais dominées par la discordance de HOD résiduelle : un changement de 2\% de $\omega_b$ impliquait des changements de plusieurs dizaines de pour cent de la 2PCF à $s\sim20\Mpch$, et les matrices de Fisher étaient singulières.

\paragraph{3. Calibration de la HOD.} La réponse à la HOD mesurée sur 72 catalogues fiduciels et soustraite des paires était prédite avec une erreur de 45\% (linéaire) et de 8\% (GP) ; avec un terme de HOD dix fois supérieur au signal, même 8\% laissent un résidu de la taille du signal.

\paragraph{4. Émulateur.} Le premier émulateur, entraîné sur 8 tirages de HOD par cosmologie, reproduisait des cosmologies exclues à une fraction de l'erreur statistique, mais ne capturait pas la réponse à $\omega_{\rm cdm}$ en dessous de $70\Mpch$ (Figure~\ref{fig:loco}). Une comparaison préliminaire avec cet émulateur donnait ASTRA et density split comparables dans l'ensemble, le density split en avance dans les vides (Q1) pour $\omega_b$ et $\omega_{\rm cdm}$, et leur combinaison 15--25\% meilleure que l'un ou l'autre.

\paragraph{5. Densification.} 240 catalogues aux HOD proches de la fiducielle ont été ajoutés dans les cosmologies de la grille de dérivées ; cela seul n'a pas stabilisé les dérivées (dispersion de 70--170\% entre réseaux).

\paragraph{6. Régression sur les paires.} Introduite comme estimateur indépendant de l'émulateur (Section~\ref{sec:regression}) ; elle a révélé que les statistiques répondent essentiellement à une combinaison des paramètres de HOD.

\paragraph{7. Entraînement sur les différences de catalogues.} A divisé par deux l'erreur des différences prédites entre catalogues exclus pour $\omega_b$ et $\sigma_8$ (21\%$\to$10\%, 13\%$\to$7\%), jusqu'à leur niveau de bruit (Tableau~\ref{tab:local}) ; le poids aux petites échelles a été abandonné car il déstabilisait les prévisions d'auto-corrélations. Avec la covariance de 192 sous-volumes, l'émulateur et la régression s'accordaient alors sur le fait qu'ASTRA était aussi bon que le density split ou meilleur de 0--20\% (Tableaux~\ref{tab:ratio} et~\ref{tab:sigma_sub}).

\paragraph{Courbe d'apprentissage.} Pour décider si davantage de catalogues d'entraînement réduiraient l'erreur de l'émulateur, nous avons réentraîné l'émulateur de la 2PCF et des corrélations croisées Q1, Q4 des quartiles aléatoires sur 25, 50, 75 et 100\% de l'ensemble d'entraînement. L'erreur de test médiane, en unités de $\sigma_{\rm box}$, passe de 0.7--1.3 à 0.5--0.8 entre 25 et 50\%, mais seulement à 0.5--0.7 à 100\% (ASTRA Q1 : 1.01, 0.59, 0.55, 0.51 ; density split Q4 : 0.69, 0.51, 0.48, 0.46). L'erreur a atteint un plateau à $\simeq0.5\sigma_{\rm box}$, comparable au bruit d'échantillonnage d'un catalogue unique : doubler l'ensemble d'entraînement ($\simeq230$ heures-nœud) l'améliorerait de $\sim10\%$ au plus, et nous ne l'avons pas fait.

\paragraph{8. Covariance.} Avec les 1786 boîtes indépendantes et les mêmes dérivées, l'avance d'ASTRA a disparu (Tableau~\ref{tab:covbox}).

\paragraph{9. Quijote.} Les prévisions ont d'abord été calculées pour les six paramètres. Un test de convergence automatique lancé pendant la nuit a rapporté que la covariance n'était pas convergée, car il comparait des prévisions complètes dans lesquelles le bruit des dérivées domine le changement ; une fois les deux effets séparés, la covariance s'est avérée convergée et le traitement supplémentaire qu'elle avait déclenché a été annulé. Les prévisions à six paramètres se sont ensuite révélées dépendre du bruit des dérivées (estimateurs standard et comprimé différant de facteurs 2--10), ce qui a conduit aux sous-ensembles de paramètres de la Section~\ref{sec:quijote}. Un premier résultat nocturne, ASTRA plus serré que le density split sur $\sigma_8$ d'un facteur 2--3 et le density split plus serré sur $\Omega_m$, $\Omega_b$, ne survit dans les sous-ensembles fiables que pour les auto-corrélations et les quartiles de halos.


\paragraph{10. MCMC.} Les premières exécutions d'inférence sur des boîtes entières étaient limitées par l'erreur de l'émulateur ; sur un sous-volume, elles sont limitées par le prior (Section~\ref{sec:mcmc}).
